% ===== 原文 pp.77--80 =====
% [原文 p.77]
% （承接 tr05c 末尾：「……特别地 $N_\ell=N$。」）
味对称群为
\begin{equation}
    \left[\prod_{k=1}^{\ell} U(n_k)\right]/U(1)_{\mathrm{diag}}.
    \tag{7.16}
\end{equation}
% [Tutor Note: p.77 (7.16) 之后的质量参数计数：由 (7.15) 下方递推 $N_j-N_{j-1}=n_j+n_{j+1}+\cdots+n_\ell$（$N_0=0$）得 $N_1=\sum_{k=1}^{\ell}n_k$；(7.16) 中 $\prod_k U(n_k)$ 的秩为 $\sum_k n_k$，商去 $U(1)_{\mathrm{diag}}$ 后剩 $\sum_k n_k-1=N_1-1$ 个质量参数，与原式一致。两个锚点可自检：满穿刺点 $Y=1^N$ 时 $n_1=N$、$\ell=1$，给 $N-1$；单穿刺点 $Y=(N-1)+1$ 时给 $1$。注意文字层会把 $n_\ell-1$ 抽成下标 $n_{\ell-1}$，必须看图像。]

该穿刺点有 $n_1+\cdots+n_\ell-1=N_1-1$ 个质量参数。幂零矩阵 (7.14) 也可以理解为提升算子（raising operator）在嵌入 $\rho\colon SU(2)\to SU(N)$ 下的像。在那样的等价描述中，味对称群 (7.16) 作为 $\rho$ 的中心化子（commutant）出现。

\para{箭图尾巴与穿刺点。} 从箭图尾巴 (7.15) 及其由 M 理论得到的 SW 解 [178] 出发，Gaiotto 在 [1] 中弄清了微分 $\phi_k$ 的极点阶数的相关模式。这帮助确认了驯服穿刺点由 $N$ 的分拆（partition）标记。带任意箭图尾巴的线性箭图在 class S 中实现为：把 $\cX(\fsu(N))$ 在带任意多个单驯服穿刺点、以及两个（一般的）驯服穿刺点的球面上约化。

事实上，我们所描述的部分 Higgsing 过程把（作为出发点的）满驯服穿刺点替换成正好由 $Y$ 标记的驯服穿刺点。正如满驯服穿刺点带有 $SU(N)$ 味对称性，该穿刺点带有味对称群 (7.16)。箭图尾巴、穿刺点与极点阶数之间的字典在 [6] 中有很好的写法。

穿刺点可以整个被闭合掉。这在 (7.8) 中的单穿刺点情形最容易看出：在规范理论一侧，相邻的两个 $SU(N)$ 群被约化成它们的对角子群，于是我们只剩下一个更短的线性箭图。

\para{部分退化的顶点算子。} 在 Toda CFT 一侧，部分 Higgsing 对应于把一个非退化顶点算子换成部分退化的顶点算子。

按脚注 56，顶点算子 $V_\alpha$ 的 1 级零 $W_N$ 后代（level-1 null $W_N$ descendants）由那些满足 $\langle\alpha-Q,e\rangle=-b-1/b$ 的根 $e$ 刻画。在 $\alpha-Q$ 的 Weyl 群变换下，我们只取这些根为单根 $e_j$。条件于是化为 $\langle\alpha,e_j\rangle=0$，即用 $\fsu(N)$ 定义表示（defining representation）的权 $h_i$ 表出为 $\langle\alpha,h_j\rangle=\langle\alpha,h_{j+1}\rangle$。$\alpha$ 的各分量 $\langle\alpha,h_i\rangle$ 组织成
\begin{equation}
    \alpha=\big(\underbrace{\alpha_{(1)},\cdots,\alpha_{(1)}}_{l_1},\cdots,\underbrace{\alpha_{(r)},\cdots,\alpha_{(r)}}_{l_r}\big)
    \tag{7.17}
\end{equation}
% [Tutor Note: p.77 (7.17) 的两个 underbrace：$l_k$ 是第 $k$ 组相等分量的个数，$r$ 是组的个数。可核对三点：(i) $l_1\ge\cdots\ge l_r\ge0$ 恰好是某个杨图 $Y$ 的列长，故这个 $Y$ 与 (7.14) 处标记幂零轨道的分拆是同一个；(ii) $\alpha$ 有 $N$ 个分量，所以 $\sum_k l_k=N$；(iii) 每个相等组只贡献一个自由值，共 $r$ 个，再减去无迹条件 $\sum_i\langle\alpha,h_i\rangle=0$ 得 $r-1$ 个参数，而 $r$ 正是 $Y$ 的列数、即 $Y$ 第一行的长度，故与原文“动量有 $r-1$ 个参数”及 (7.16) 的计数都吻合。]

其中 $l_k$，$1\le k\le r$，表示相等分量的个数，并且我们可以重新排列这些分量，使得 $l_1\ge l_2\ge\cdots\ge l_r\ge0$ 给出某个杨图 $Y$ 的列长。这与穿刺点的杨图分类是同一个分类。例如动量有 $r-1$ 个参数（因为无迹性），亦即 $Y$ 第一行的长度（减一），从而与 (7.16) 中的计数相符。

% [原文 p.78]
[336] 中的精确提法是：由 $Y$ 标记的驯服穿刺点对应于动量为
\begin{equation}
    \alpha=\left(b+\frac{1}{b}\right)\rho_Y+m_Y,
    \tag{7.18}
\end{equation}
的顶点算子，其中 $m_Y$ 是味对称群 (7.16) 的质量参数，$\rho_Y$ 是 Weyl 向量 $\rho$ 到具有重数 (7.17) 的子空间上的投影。对满穿刺点情形 $Y=1^N$ 有 $\rho_Y=\rho$，而对单穿刺点 $Y=(N-1)+1$ 可得 $\rho_Y=(N/2)\varpi_1$。两种情形下这个提法都重现 (7.9)。
% [Tutor Note: p.78 (7.18) 后的两个特例可手算核对：$\fsu(N)$ 的 Weyl 向量按定义表示的权写出为 $\rho=\mathrm{diag}\big(\frac{N-1}{2},\frac{N-3}{2},\ldots,\frac{1-N}{2}\big)$。(i) 满穿刺点 $Y=1^N$（各列长都是 1，无重数可平均）→ $\rho_Y=\rho$。(ii) 单穿刺点 $Y=(N-1)+1$（一列长 $N-1$、一列长 1）：让 $N-1$ 个较小分量成块，其和 $=-\frac{N-1}{2}$、平均 $=-\frac12$，剩下最大的分量 $\frac{N-1}{2}$ 单独成块，于是 $\rho_Y=\big(-\frac12,\ldots,-\frac12,\frac{N-1}{2}\big)=\frac{N}{2}\varpi_1$，其中 $\varpi_1=\mathrm{diag}(1,0,\ldots,0)-\frac1N\id$。注意若反过来把最大的 $N-1$ 个分量成块，会得到 $\big(\frac12,\ldots,\frac12,\frac{1-N}{2}\big)$，与原文不符；故 (7.17) 的分块必须按「重数递减、取值递增」的约定理解。]

\para{穿刺点：幂零轨道、Nahm、Hitchin 与 Toda。} 对 A 型情形 $\fg=\fsu(N)$，我们已经看到驯服穿刺点由 $N$ 的分拆标记。更宽广的视角是：$\cX(\fsu(N))$ 的驯服余维 2 缺陷由这样的分拆标记。D 型情形以及 USp--SO 箭图尾巴在 [6] 中讨论。

对一般的 6d (2,0) 理论 $\cX(\fg)$，有（至少）三套数据可以等价地刻画该缺陷。\footnote{欢迎帮助补全参考文献。}
\begin{itemize}[leftmargin=2em,itemsep=2pt,topsep=2pt]
  \item \textbf{Nahm 数据。} 一个幂零轨道 $\cO_N\subset\fg_\C$，它描述规范群为 $G$ 的 SYM 的 Nahm 极点边界条件。这来自于把 $\cX(\fg)$ 放在 $\R^{2,1}\times\mathrm{cigar}\times S^1$ 上、并令缺陷位于 cigar 的尖端：先沿 cigar 的圆周方向约化，给出规范群为 $G$ 的 5d $\cN=2$ SYM，并带 Nahm 极点边界条件。
  \item \textbf{Hitchin 数据。} 一个幂零轨道 $\cO_H\subset\fg_\C$ 外加一些离散数据。在同一个 $\R^{2,1}\times\mathrm{cigar}\times S^1$ 几何设置里，先沿 $S^1$ 约化给出 5d 中的一个余维 2 缺陷，再把 cigar 约化为半直线，就给出 Nahm 极点边界条件的 S 对偶。这套数据很早就有人在 [356] 中研究过。
  \item \textbf{Toda 数据。} $\fg$ Toda CFT 的一个部分退化主算子，由其零 $W_\fg$ 后代指定。
\end{itemize}
% [Tutor Note: p.78 三套数据为什么等价：缺陷放在 $\R^{2,1}\times\mathrm{cigar}\times S^1$ 的 cigar 尖端，两条圆周（cigar 的 $S^1$ 纤维圆与显式的 $S^1$）谁先约化给出两条不同的 5d 描述。先收 cigar 的纤维圆 → 5d $\cN=2$ SYM 加 Nahm 极点边界条件（$e$ 幂零才不破坏标度不变性，这是取“幂零轨道”而非任意伴随轨道的来由）；先收显式的 $S^1$ → 5d 的余维 2 缺陷，再把 cigar 缩成半直线 → 该边界条件的 S 对偶，即 Hitchin 侧数据。两种约化顺序可交换，正是三分类等价的物理基础；Toda 侧则是同一个 $\fg$ 的 W 代数退化表示，故 $\cO_N$、$\cO_H$ 与零 $W_\fg$ 后代彼此对应。]

Nahm 数据与 Hitchin 数据在 [100] 中被联系起来。与 Toda 数据的关系则在 [318] 中对 $e$ 是 $\fg$ 的某个 Levi 子代数内的主幂零元（principal nilpotent）的情形得到理解，另见 [357,358]。

同一分类对李代数为 $\fg$ 的 6d (2,0) 小弦理论中的余维 2 算子也成立（这蕴含了 $\cX(\fg)$ 的分类）。[342,359,360] 中的做法是把小弦实现为 $\C^2/\Gamma$ 奇点上的 IIB 弦，其中缺陷由包裹 $\C^2/\Gamma$ 中某些非紧二维闭链（non-compact two-cycles）的 D5 膜组成。

\para{非单链的规范群。} 把 $\cX(\fg)$ 在圆周上约化给出规范代数 $\fg$ 的 5d $\cN=2$ SYM，而 $\fg$ 是单链的。非单链的规范群通过扭（twisting）来实现，也就是用 $\fg$ 的一个外自同构改变场的周期性。\footnote{这里“扭”的概念与把 $\cX(\fg)$ 在 $C$ 上约化时为保持超对称所用的部分拓扑扭曲无关。}

% [原文 p.79]
\begin{figure}[htbp]
\centering
\includegraphics[width=0.95\textwidth]{fig4.png}
\figtitle{4}{左：由交替的 USp$(2N-2)$ 与 Spin$(2N)$ 群构成的线性箭图，以及其带分支割线的 $\fso(2N)$ class S 曲线。右：对 $N=2$，同一理论作为 $SU(2)$ 广义箭图，以及其 $\fsu(2)$ class S 曲线；该曲线是左侧 $\fso(4)$ 曲线的双重覆盖。}
\end{figure}

同样地，class S 也包含具有任意规范群的 4d $\cN=2$ 理论：它们通过引入（外自同构）扭线（twist lines）得到，而扭线可以终止于穿刺点。这一方向在 [96,98,99,318] 中被探索（另见 [361]）。$\fg=\fso(2N)$ 的例子、以及它在 $N=2$ 时因 $\fso(4)=\fsu(2)^2$ 而重新表述为 $\fsu(2)$ class S 理论的例子，见 Figure 4。

\subsection{野生穿刺点与 AD 理论}

我们迄今的考察只涉及 6d (2,0) 理论中所谓的“驯服”余维 2 缺陷。它们可以广泛推广到“野生”缺陷，后者由 Witten [362] 在 4d $\cN=4$ SYM 的面算子语境中引入。这些缺陷在其支撑附近对 6d 理论的“场”施加更强的爆破（blow-up）。

\para{来自碰撞的野生穿刺点。} 我们回忆有质量的驯服穿刺点 (4.2)
\begin{equation}
    \varphi(z)\sim\left(\frac{\diag(m,-m)}{z-z_i}+O(1)\right)\dd z
    \quad\Longrightarrow\quad
    \phi_2(z)=\left(\frac{m^2}{(z-z_i)^2}+O\left(\frac{1}{z-z_i}\right)\right)\dd z^2
    \tag{7.19}
\end{equation}
及其无质量版本 (4.4)。让 $\varphi$ 的 $l$ 个这样的单极点相撞、同时适当地标度质量参数，就导致 $l$ 阶极点，于是一般地给出 $\phi_2\sim\dd z^2/(z-z_i)^{2l}$。正如其驯服对应物那样，所得的 $l$ 级野生穿刺点可以通过对 $\Phi_z$ 级数展开中的本征值施加某些关系而被部分闭合，使得 $\phi_2$ 的极点阶数低于 $2l$。
% [Tutor Note: p.79 (7.19) 可直接核出碰撞极限的标度律：要得到 $\phi_2\sim\dd z^2/(z-z_i)^{2l}$，由 $m^2/(z-z_i)^2$ 需 $m^2\sim(z-z_i)^{-2l+2}$，即 $m\sim\epsilon^{-(l-1)}$（$\epsilon=z-z_i\to0$）。$l=2$ 时 $m\sim1/\epsilon$，这正是原文“适当地标度质量参数”的内容；若质量固定不动，两个单极点相撞后 $\phi_2$ 仍只有 2 阶极点（即驯服穿刺点）。另外注意 (7.19) 左端用 $\varphi$（标量场）、右端用 $\phi_2$（二次微分系数），全文须统一：$\varphi$ 是 4d 标量，$\phi_k=u_k\,\dd z^k$ 是 SW 微分，$\Phi_z$ 是带本征值的微分矩阵。]

碰撞极限对 4d 规范理论可以有两大类效应：通过让一些超多重态变重而退耦它们、同时保持动力学标度 $\Lambda$ 固定；或者把理论调到 Coulomb 分支上的一个 AD 点 [2,363--365]，在该点理论变成强耦合的孤立 SCFT。

对眼下我们所考虑的 $\fg=\fsu(2)$ 情形，野生穿刺点由 $\phi_2$ 的极点阶数来标记（对驯服穿刺点该阶数为 2），当然也由这些极点处展开的系数标记。像驯服情形那样沿圆周切割黎曼面，$\fsu(2)$ class S 理论就可以通过规范化 $SU(2)$ 味对称性来构造

% [原文 p.80]
\begin{figure}[htbp]
\centering
\includegraphics[width=0.7\textwidth]{fig5.png}
\figtitle{5}{把典型的 $\fsu(2)$ class S 理论（对应带两个驯服穿刺点与两个野生穿刺点的球面）分解为 $T_2$ 与 $X_p$ 两种构件，它们通过规范化对角的 $\fsu(2)$ 味对称性耦合在一起。}
\end{figure}

三基本半超多重态（trifundamental half-hypermultiplet）$T_2$（对应带三个驯服穿刺点的球面），以及理论 $X_p$（对应带一个驯服穿刺点和一个 $\phi_2$ 具有 $p>2$ 阶极点的野生穿刺点的球面）。见 Figure 5。带单个野生穿刺点的球面无法切分成这些构件，并导致其他有趣的理论 $Y_p$。这就穷尽了 $\fsu(2)$ class S。

\para{带野生穿刺点的理论示例。} 只为这一解释，我们用 $(p_1\ p_2\ \ldots\ p_k)$ 记由带 $k$ 个穿刺点、且各点处 $\phi_2$ 分别有 $p_1,\ldots,p_k$ 阶极点的球面所得到的 class S 理论。让我们用这个记法下的 $(2\ 2\ 2\ 2)$------即取 $C$ 为带四个驯服穿刺点的球面所实现的 $SU(2)$ $N_f=4$ SQCD------出发，来例示上面两类效应。我们先退耦超多重态。
\begin{itemize}[leftmargin=2em,itemsep=2pt,topsep=2pt]
  \item $SU(2)$ $N_f=3$ SQCD 来自 $(2\ 2\ 4)$，即带两个驯服穿刺点和一个 4 阶野生穿刺点的球面，该野生穿刺点由两个驯服穿刺点相撞得到。
  \item $SU(2)$ $N_f=2$ SQCD 在 class S 中以两种方式出现。第一，作为由 $(2\ 2\ 4)$ 让两个驯服穿刺点相撞所得的 $(4\ 4)$。另一种方式是作为 $(2\ 2\ 3)$：可以调节 $N_f=3$ 理论的 $(2\ 2\ 4)$ 描述中野生穿刺点的某个质量参数来退耦该超多重态，这会把该野生穿刺点处 $\phi_2$ 的极点从 4 阶降低到 3 阶。一个一致性检验是：两种构造给出等价的 SW 几何。
  \item $SU(2)$ $N_f=1$ SQCD 于是作为 $(4\ 3)$ 出现。
  \item 纯 $SU(2)$ SYM 作为 $(3\ 3)$ 出现，带两个最小野生穿刺点。
\end{itemize}

还有进一步的碰撞极限，结果它们实现的是 AD 理论。让 $SU(2)$ $N_f=1$ SQCD 的 $(4\ 3)$ 实现中的两个野生穿刺点相撞，我们得到一个阶数相当高的单个野生穿刺点 (7)：这就是上面提到的 $Y_7$ 理论。$SU(2)$ $N_f=2$ 的 Coulomb 分支的 AD 点（最奇异的点）通过让穿刺点 $(4\ 4)\to(8)$ 或 $(2\ 2\ 3)\to(2\ 5)$ 相撞得到，两条带穿刺点的曲线 $C$ 最终给出同一个 4d SCFT（即 $Y_8\simeq X_5$）。对 $SU(2)$ $N_f=3$ 我们找到碰撞 $(2\ 2\ 4)\to(2\ 6)$，即 $X_6$。当然，这些极限在规范理论一侧都翻译成调节参数，因而很久以前就已被发现 [2,363]，但 class S 的实现把它们嵌入了更宽广的框架之中。
% [Tutor Note: p.80 用 $(p_1\ \cdots\ p_k)$ 记法做阶数守恒检验（两个穿刺点相撞时 $\phi_2$ 极点阶数相加）：$(2\ 2\ 2\ 2)\to(2\ 2\ 4)$（$2+2=4$）、$(2\ 2\ 4)\to(4\ 4)$（$2+2=4$）、$(4\ 3)\to(7)$（$4+3=7$）、$(4\ 4)\to(8)$（$4+4=8$）、$(2\ 2\ 3)\to(2\ 5)$（$2+3=5$）、$(2\ 2\ 4)\to(2\ 6)$（$2+4=6$），全部吻合。但 $(2\ 2\ 4)\to(2\ 2\ 3)$ 不是相撞、而是把野生穿刺点“部分闭合”（调节质量参数），故得到 3 阶而非 4 阶------这正是原文说 $N_f=2$ 有两种 class S 实现的原因。$X_p$ 与 $Y_p$ 的名字也由此记法来：$X_5=(2\ 5)$ 是单穿刺点加 5 阶野生点，$Y_8$ 是只带一个 8 阶野生点的球面，故 $Y_8\simeq X_5$ 是同一 SCFT 的两种描述。]
